\section{Displayed data I --- a Tate-style conductor numeral at
$(A,B)=(29,1)$}\label{sec:tate}

The file
\texttt{Tate\_Frey\_Conductor\_29\_Neron\_final.lean}
(\texttt{11c0ad7}; foundations incremental
\href{https://github.com/DavidFox998/beal-level-26-foundations/actions/runs/35488239286}{35488239286},
$137$s cache hit) packages the displayed conductor numeral, for the
single specialization $(A,B)=(29,1)$ of the gap-3 equation
\eqref{eq:gap3},
\[
  N_E^{\mathrm{disp}}=(2:\mathbb{N})^5\cdot 29=928.
\]
The named theorems \texttt{displayed\_Neron\_conductor} and
\texttt{final\_N\_E\_eq\_928} are that equality, proved by
\texttt{decide}. The split \texttt{displayed\_928\_div\_29} is
$928/29=32$. At the other named specialization $(A,B)=(1,1)$ the
same family displays conductor $32$, not $928$
(\texttt{sagemath/certs/tate\_nero\_29.json}): the numeral $928$ is
a fact about one pair $(A,B)$, not a property of every solution of
\eqref{eq:gap3}.

\begin{displaythm}[Displayed N\'eron conductor]
\label{disp:tate-928}
In Lean~4.12, with \texttt{open Nat Finset Classical}:
\begin{itemize}[leftmargin=1.6em]
\item $(2:\mathbb{N})^5\cdot 29=928$ (\texttt{final\_N\_E\_eq\_928});
\item $928/29=32$ (\texttt{displayed\_928\_div\_29});
\item $32\cdot 29=928$;
\item $\neg\,29\mid 32$.
\end{itemize}
These are numerals. Mathlib~$4.12$ has no \texttt{KodairaType} and
no \texttt{NeronModel}. The name $N_E=928$ is not the output of a
Lean formalization of Tate's algorithm or of a N\'eron model.
\end{displaythm}

The displayed Weierstrass invariants of~\eqref{eq:frey} are
\[
  c_4=16(A^8+B^8+A^4B^4),\qquad
  \Delta=16A^8B^8(A^4+B^4)^2.
\]
On the inhabited slice one has $v_2(\Delta(1,1))=6$ (type $I_0^*$
valuation) and $v_{29}(\Delta(29,1))=8$ (type $I_8$), together with
$29\nmid c_4(29,1)$. The residual-level Finset $\{32,928\}$ has
card~$2$; that card is a Finset numeral over two named
specializations, not two N\'eron models of a single curve. The
parent names \texttt{Tate\_algorithm\_at\_29},
\texttt{Frey\_Neron\_conductor}, and
\texttt{Frey\_conductor\_29\_is\_Neron} are \emph{theorems in this
namespace}: each re-exports the inhabited conjunction
\texttt{Tate\_algorithm\_at\_2\_and\_29}. On the parent file
\texttt{Tate\_Frey\_Conductor\_29.lean} those three names stay
\texttt{def Prop}.

\begin{honesty}
The phrase ``final $N_E=928$'' means the displayed identity
$2^5\cdot 29=928$ at $(A,B)=(29,1)$. It is not Tate's algorithm as a
function on integral models, and it is not a claim about the
conductor for other $(A,B)$ satisfying \eqref{eq:gap3}. SAGE
\texttt{sagemath/tate\_nero\_29.sage} uses the displayed
Frey~\eqref{eq:frey}, not the short model, and its own comment
records that $(1,1)$ gives $32$.
\end{honesty}

Axioms on the packaged Tate final are $[\mathtt{propext}]$ for the
pure numerals and $[\mathtt{propext},\mathtt{Classical.choice},\mathtt{Quot.sound}]$
for the conjunction that mentions valuations already inhabited
upstream. There is no \texttt{Lean.ofReduceBool}, no
\texttt{sorryAx}, and no new axiom.
